
\subsubsection{4.1 Research Questions}

\textbf{RQ1:} Does DNSI correlate with known generalization gaps across benchmarks with ground truth held-out evaluations?

\textbf{RQ2:} Can pre-saturation DNSI values predict future generalization gaps?

\textbf{RQ3:} Does the DNSI-gap correlation generalize across modalities (vision and NLP)?

\subsubsection{4.2 Datasets}

Evaluation uses four benchmarks with published ground truth generalization gap measurements:

\begin{table}[htbp]
\centering
\resizebox{\columnwidth}{!}{%
\begin{tabular}{llll}
\toprule
Benchmark & Domain & Gap Source & Ground Truth Gap \\
\midrule
ImageNet & Vision & Recht et al. (2019) & 12.5\% \\
CIFAR-10 & Vision & Recht et al. (2019) & 4.0\% \\
ObjectNet & Vision & Barbu et al. (2019) & 42.5\% \\
HANS & NLP & McCoy et al. (2019) & 40.0\% \\
\bottomrule
\end{tabular}
}
\end{table}

\subsubsection{4.3 Baselines}

DNSI is compared against: (1) Raw Entropy (no difficulty normalization), (2) Improvement Rate (mean accuracy gain per year), and (3) Time Since Last Improvement.

\subsubsection{4.4 Implementation Details}

\begin{itemize}
\item Window size: 6 months
\item Minimum SOTA entries: 15
\item Bootstrap samples: 10,000
\item Random seed: 42
\end{itemize}

\subsubsection{4.5 Evaluation Metrics}

\textbf{Primary:} Pearson R between DNSI and generalization gap (threshold: R < -0.4)

\textbf{Secondary:} Spearman $\rho$, bootstrap confidence intervals, $R^2$ for temporal prediction

\subsubsection{4.6 Limitations of Experimental Setup}

The proof-of-concept uses synthetic SOTA histories that preserve realistic saturation dynamics but are not drawn directly from PapersWithCode data. DNSI values for ImageNet, ObjectNet, and HANS are synthetic estimates based on the methodology validated in CIFAR-10 analysis. NLP benchmarks lack a universal difficulty proxy equivalent to class count.
